% Leaf5 版 — 每页 2 张幻灯片（1650x1264屏幕）


% ====== 引言（9 页） ======

\leafslidepair{1引言}{5}{39}
\leafslidepair{1引言}{40}{41}
\leafslidepair{1引言}{42}{43}
\leafslidepair{1引言}{44}{45}
\leafslidepair{1引言}{46}{0}

% ── 引言 章末笔记 ──


% ====== 数学基础（13 页） ======

\leafslidepair{2数学基础}{3}{4}
\leafslidepair{2数学基础}{5}{6}
\leafslidepair{2数学基础}{7}{8}
\leafslidepair{2数学基础}{9}{10}
\leafslidepair{2数学基础}{11}{12}
\leafslidepair{2数学基础}{13}{14}
\leafslidepair{2数学基础}{15}{0}

% ── 数学基础 章末笔记 ──


% ====== 凸分析（28 页） ======

\leafslidepair{3凸分析}{3}{4}
\leafslidepair{3凸分析}{5}{8}
\leafslidepair{3凸分析}{9}{11}
\leafslidepair{3凸分析}{12}{14}
\leafslidepair{3凸分析}{15}{17}
\leafslidepair{3凸分析}{18}{21}
\leafslidepair{3凸分析}{23}{24}
\leafslidepair{3凸分析}{25}{27}
\leafslidepair{3凸分析}{29}{30}
\leafslidepair{3凸分析}{34}{39}
\leafslidepair{3凸分析}{40}{41}
\leafslidepair{3凸分析}{42}{43}
\leafslidepair{3凸分析}{44}{45}
\leafslidepair{3凸分析}{46}{47}

% ── 凸分析 章末笔记 ──


% ====== 线性规划基本性质（26 页） ======

\leafslidepair{4线性规划基本性质}{3}{4}
\leafslidepair{4线性规划基本性质}{5}{6}
\leafslidepair{4线性规划基本性质}{7}{8}
\leafslidepair{4线性规划基本性质}{9}{10}
\leafslidepair{4线性规划基本性质}{11}{12}
\leafslidepair{4线性规划基本性质}{13}{14}
\leafslidepair{4线性规划基本性质}{15}{16}
\leafslidepair{4线性规划基本性质}{17}{21}
\leafslidepair{4线性规划基本性质}{22}{23}
\leafslidepair{4线性规划基本性质}{24}{25}
\leafslidepair{4线性规划基本性质}{26}{27}
\leafslidepair{4线性规划基本性质}{28}{29}
\leafslidepair{4线性规划基本性质}{30}{31}

% ── 线性规划基本性质 章末笔记 ──


% ====== 单纯形方法（40 页） ======

\leafslidepair{5单纯形方法}{2}{3}
\leafslidepair{5单纯形方法}{4}{5}
\leafslidepair{5单纯形方法}{6}{7}
\leafslidepair{5单纯形方法}{12}{15}
\leafslidepair{5单纯形方法}{16}{24}
\leafslidepair{5单纯形方法}{25}{26}
\leafslidepair{5单纯形方法}{27}{28}
\leafslidepair{5单纯形方法}{29}{30}
\leafslidepair{5单纯形方法}{35}{36}
\leafslidepair{5单纯形方法}{37}{38}
\leafslidepair{5单纯形方法}{39}{40}
\leafslidepair{5单纯形方法}{41}{42}
\leafslidepair{5单纯形方法}{43}{44}
\leafslidepair{5单纯形方法}{45}{46}
\leafslidepair{5单纯形方法}{47}{48}
\leafslidepair{5单纯形方法}{49}{50}
\leafslidepair{5单纯形方法}{51}{52}
\leafslidepair{5单纯形方法}{53}{54}
\leafslidepair{5单纯形方法}{55}{56}
\leafslidepair{5单纯形方法}{57}{58}

% ── 单纯形方法 章末笔记 ──


% ====== 对偶理论（57 页） ======

\leafslidepair{6对偶理论}{3}{4}
\leafslidepair{6对偶理论}{5}{6}
\leafslidepair{6对偶理论}{7}{8}
\leafslidepair{6对偶理论}{9}{10}
\leafslidepair{6对偶理论}{11}{15}
\leafslidepair{6对偶理论}{16}{17}
\leafslidepair{6对偶理论}{18}{19}
\leafslidepair{6对偶理论}{21}{22}
\leafslidepair{6对偶理论}{23}{24}
\leafslidepair{6对偶理论}{25}{26}
\leafslidepair{6对偶理论}{27}{28}
\leafslidepair{6对偶理论}{29}{30}
\leafslidepair{6对偶理论}{31}{32}
\leafslidepair{6对偶理论}{33}{40}
\leafslidepair{6对偶理论}{41}{42}
\leafslidepair{6对偶理论}{43}{44}
\leafslidepair{6对偶理论}{45}{46}
\leafslidepair{6对偶理论}{47}{48}
\leafslidepair{6对偶理论}{49}{50}
\leafslidepair{6对偶理论}{51}{52}
\leafslidepair{6对偶理论}{53}{54}
\leafslidepair{6对偶理论}{55}{56}
\leafslidepair{6对偶理论}{57}{58}
\leafslidepair{6对偶理论}{59}{60}
\leafslidepair{6对偶理论}{61}{62}
\leafslidepair{6对偶理论}{63}{64}
\leafslidepair{6对偶理论}{65}{66}
\leafslidepair{6对偶理论}{67}{68}
\leafslidepair{6对偶理论}{81}{0}

% ── 对偶理论 章末笔记 ──


% ====== 最优性条件（43 页） ======

\leafslidepair{7最优性条件}{3}{4}
\leafslidepair{7最优性条件}{5}{7}
\leafslidepair{7最优性条件}{13}{14}
\leafslidepair{7最优性条件}{15}{16}
\leafslidepair{7最优性条件}{17}{18}
\leafslidepair{7最优性条件}{19}{20}
\leafslidepair{7最优性条件}{21}{22}
\leafslidepair{7最优性条件}{23}{24}
\leafslidepair{7最优性条件}{25}{26}
\leafslidepair{7最优性条件}{27}{28}
\leafslidepair{7最优性条件}{29}{34}
\leafslidepair{7最优性条件}{35}{36}
\leafslidepair{7最优性条件}{37}{38}
\leafslidepair{7最优性条件}{39}{40}
\leafslidepair{7最优性条件}{41}{42}
\leafslidepair{7最优性条件}{43}{48}
\leafslidepair{7最优性条件}{49}{50}
\leafslidepair{7最优性条件}{51}{52}
\leafslidepair{7最优性条件}{53}{54}
\leafslidepair{7最优性条件}{55}{56}
\leafslidepair{7最优性条件}{57}{58}
\leafslidepair{7最优性条件}{59}{0}

\leafnotepage{
  \section*{FJ vs KKT 条件对比}

  \textbf{FJ条件}：存在不全为零的 $\lambda_0,\dots,\lambda_m,\mu_1,\dots,\mu_l$ 使得
  $\lambda_0\nabla f + \sum\lambda_i\nabla g_i + \sum\mu_j\nabla h_j = 0$，
  $\lambda_i \ge 0$，$\lambda_i g_i = 0$。
  注意：$\lambda_0$ 可以为 0，此时不含目标函数信息。

  \textbf{KKT条件}：增加约束规格后，保证 $\lambda_0 \neq 0$（可归一化为1）。
  $\nabla f + \sum\lambda_i\nabla g_i + \sum\mu_j\nabla h_j = 0$，
  $\lambda_i \ge 0$，$\lambda_i g_i = 0$。
  即 Lagrange 函数 (x,\lambda,\mu)$ 的驻点条件。

  真题链接：2025 Q6 — 先求所有FJ点，再证不存在KKT点。
}

% ====== 算法（19 页） ======

\leafslidepair{8算法}{3}{4}
\leafslidepair{8算法}{5}{6}
\leafslidepair{8算法}{7}{8}
\leafslidepair{8算法}{9}{10}
\leafslidepair{8算法}{14}{15}
\leafslidepair{8算法}{16}{17}
\leafslidepair{8算法}{18}{19}
\leafslidepair{8算法}{20}{21}
\leafslidepair{8算法}{22}{23}
\leafslidepair{8算法}{24}{0}

% ── 算法 章末笔记 ──


% ====== 一维搜索（17 页） ======

\leafslidepair{9一维搜索}{3}{4}
\leafslidepair{9一维搜索}{5}{6}
\leafslidepair{9一维搜索}{7}{26}
\leafslidepair{9一维搜索}{27}{51}
\leafslidepair{9一维搜索}{52}{53}
\leafslidepair{9一维搜索}{54}{55}
\leafslidepair{9一维搜索}{56}{57}
\leafslidepair{9一维搜索}{58}{59}
\leafslidepair{9一维搜索}{60}{0}

% ── 一维搜索 章末笔记 ──


% ====== 使用导数的优化算法（79 页） ======

\leafslidepair{10使用导数的优化算法}{40}{41}
\leafslidepair{10使用导数的优化算法}{42}{43}
\leafslidepair{10使用导数的优化算法}{44}{45}
\leafslidepair{10使用导数的优化算法}{46}{47}
\leafslidepair{10使用导数的优化算法}{48}{49}
\leafslidepair{10使用导数的优化算法}{50}{51}
\leafslidepair{10使用导数的优化算法}{52}{53}
\leafslidepair{10使用导数的优化算法}{54}{55}
\leafslidepair{10使用导数的优化算法}{56}{57}
\leafslidepair{10使用导数的优化算法}{58}{59}
\leafslidepair{10使用导数的优化算法}{60}{61}
\leafslidepair{10使用导数的优化算法}{62}{63}
\leafslidepair{10使用导数的优化算法}{64}{65}
\leafslidepair{10使用导数的优化算法}{66}{67}
\leafslidepair{10使用导数的优化算法}{68}{69}
\leafslidepair{10使用导数的优化算法}{70}{71}
\leafslidepair{10使用导数的优化算法}{72}{73}
\leafslidepair{10使用导数的优化算法}{74}{75}
\leafslidepair{10使用导数的优化算法}{76}{77}
\leafslidepair{10使用导数的优化算法}{78}{79}
\leafslidepair{10使用导数的优化算法}{80}{81}
\leafslidepair{10使用导数的优化算法}{82}{83}
\leafslidepair{10使用导数的优化算法}{84}{85}
\leafslidepair{10使用导数的优化算法}{86}{87}
\leafslidepair{10使用导数的优化算法}{88}{89}
\leafslidepair{10使用导数的优化算法}{90}{91}
\leafslidepair{10使用导数的优化算法}{92}{93}
\leafslidepair{10使用导数的优化算法}{94}{95}
\leafslidepair{10使用导数的优化算法}{96}{97}
\leafslidepair{10使用导数的优化算法}{98}{99}
\leafslidepair{10使用导数的优化算法}{100}{101}
\leafslidepair{10使用导数的优化算法}{102}{103}
\leafslidepair{10使用导数的优化算法}{104}{105}
\leafslidepair{10使用导数的优化算法}{106}{107}
\leafslidepair{10使用导数的优化算法}{108}{109}
\leafslidepair{10使用导数的优化算法}{110}{111}
\leafslidepair{10使用导数的优化算法}{112}{113}
\leafslidepair{10使用导数的优化算法}{114}{115}
\leafslidepair{10使用导数的优化算法}{116}{117}
\leafslidepair{10使用导数的优化算法}{118}{0}

\leafnotepage{
  \section*{拟牛顿法公式速查}

  \textbf{DFP（校Hessian逆$）}：
   = x_{k+1}-x_k,\ y_k = g_{k+1}-g_k$
  {k+1} = H_k + \frac{s_k s_k^T}{s_k^T y_k} - \frac{H_k y_k y_k^T H_k}{y_k^T H_k y_k}$
  搜索方向： = -H_k\nabla f(x_k)$

  \textbf{BFGS（校Hessian $）}：
  {k+1} = B_k + \frac{y_k y_k^T}{y_k^T s_k} - \frac{B_k s_k s_k^T B_k}{s_k^T B_k s_k}$
  搜索方向：解  d_k = -\nabla f(x_k)$

  真题链接：2025 Q4 — DFP法完整两轮迭代。
}

% ====== 无约束优化直接方法（0 页） ======

% （本章不考，已全部排除）


% ====== 可行方向法（60 页） ======

\leafslidepair{12可行方向法}{2}{3}
\leafslidepair{12可行方向法}{4}{5}
\leafslidepair{12可行方向法}{6}{7}
\leafslidepair{12可行方向法}{8}{9}
\leafslidepair{12可行方向法}{10}{11}
\leafslidepair{12可行方向法}{12}{13}
\leafslidepair{12可行方向法}{14}{15}
\leafslidepair{12可行方向法}{16}{17}
\leafslidepair{12可行方向法}{18}{19}
\leafslidepair{12可行方向法}{20}{21}
\leafslidepair{12可行方向法}{22}{23}
\leafslidepair{12可行方向法}{24}{25}
\leafslidepair{12可行方向法}{26}{27}
\leafslidepair{12可行方向法}{28}{29}
\leafslidepair{12可行方向法}{30}{31}
\leafslidepair{12可行方向法}{32}{33}
\leafslidepair{12可行方向法}{39}{40}
\leafslidepair{12可行方向法}{41}{42}
\leafslidepair{12可行方向法}{43}{44}
\leafslidepair{12可行方向法}{45}{46}
\leafslidepair{12可行方向法}{47}{48}
\leafslidepair{12可行方向法}{49}{50}
\leafslidepair{12可行方向法}{51}{52}
\leafslidepair{12可行方向法}{53}{54}
\leafslidepair{12可行方向法}{55}{56}
\leafslidepair{12可行方向法}{57}{58}
\leafslidepair{12可行方向法}{59}{60}
\leafslidepair{12可行方向法}{61}{62}
\leafslidepair{12可行方向法}{63}{64}
\leafslidepair{12可行方向法}{65}{66}

\leafnotepage{
  \section*{Rosen梯度投影法}

  投影矩阵： = I - M^T(MM^T)^{-1}M$，方向： = -P\,\nabla f(x)$

  \textbf{步骤}：
  1. 识别积极约束 → 构造矩阵 $
  2.  = I - M^T(MM^T)^{-1}M$
  3.  = -P\nabla f(x)$
  4. 若  \neq 0$：一维搜索；若 =0$：计算
      = -(MM^T)^{-1}M\nabla f(x)$
     所有  \ge 0 \to$ KKT点
     有负分量 $\to$ 去掉对应约束行重算

  \textbf{Zoutendijk法}：解LP子问题 $\min \nabla f^T d$，
  s.t. $\nabla g_i^T d + g_i \le 0$，hBc1 \le d_j \le 1$
  最优值$<0$得下降方向；$=0$为FJ/KKT点。

  真题链接：2025 Q5。
}

% ====== 罚函数法（48 页） ======

\leafslidepair{13罚函数法}{2}{3}
\leafslidepair{13罚函数法}{4}{5}
\leafslidepair{13罚函数法}{6}{7}
\leafslidepair{13罚函数法}{8}{9}
\leafslidepair{13罚函数法}{10}{11}
\leafslidepair{13罚函数法}{12}{13}
\leafslidepair{13罚函数法}{14}{15}
\leafslidepair{13罚函数法}{16}{17}
\leafslidepair{13罚函数法}{18}{19}
\leafslidepair{13罚函数法}{20}{21}
\leafslidepair{13罚函数法}{22}{23}
\leafslidepair{13罚函数法}{24}{25}
\leafslidepair{13罚函数法}{26}{27}
\leafslidepair{13罚函数法}{28}{29}
\leafslidepair{13罚函数法}{30}{31}
\leafslidepair{13罚函数法}{32}{33}
\leafslidepair{13罚函数法}{34}{35}
\leafslidepair{13罚函数法}{36}{37}
\leafslidepair{13罚函数法}{38}{39}
\leafslidepair{13罚函数法}{40}{41}
\leafslidepair{13罚函数法}{42}{43}
\leafslidepair{13罚函数法}{44}{45}
\leafslidepair{13罚函数法}{46}{47}
\leafslidepair{13罚函数法}{48}{49}

% ── 罚函数法 章末笔记 ──

